% !TeX root = ../main.tex
\section{Introduction}
\label{sec:introduction}

Translating requirements into architectural decisions requires more than
identifying components that implement requested functionality. Architects must
interpret the operational context, consider quality attributes, and justify
choices whose consequences affect subsequent system evolution~\cite{jansen2005decisions,kruchten2006architectural}.
For AI-enabled systems, this task also involves concerns spanning data
preparation, model training, deployment, and monitoring~\cite{washizaki2019patterns,heiland2023patterns}.
A user story and its acceptance criteria may describe expected behavior without
specifying the data and model lifecycle assumptions needed to select an
architectural solution. Decision support must therefore distinguish what the
requirements establish from what would require additional information.

Architectural patterns provide a vocabulary of reusable solutions to recurring
problems. Catalogs for AI-based systems organize this knowledge and make
candidate solutions accessible~\cite{heiland2023patterns,aiPatternsRepository}.
However, the existence of a relevant pattern does not establish its applicability
to a particular requirement. Selection requires examining the pattern's problem,
context, and consequences against the available requirement evidence, rather
than relying on shared terminology alone. Reuse also needs an explicit rationale:
an Architectural Decision Record (ADR) can document the context, proposed
decision, and consequences so that the choice remains inspectable and can be
revisited~\cite{nygard2011adr}. Pattern identification and decision documentation
are thus connected but distinct architectural activities.

Large language models (LLMs) offer opportunities to assist these activities by
interpreting natural-language inputs and drafting architectural
artifacts~\cite{cervantes2026assistants,dhar2024generating}.
Nevertheless, fluent output is not evidence of a sound architectural decision.
Empirical research on LLM architectural knowledge identifies limitations in
precision and the recovery of design rationale~\cite{soliman2025architectural}.
A plausible recommendation may depend on assumptions absent from the
requirements, while a well-structured ADR may still contain unsupported
reasoning. Neither a model-generated applicability judgment nor successful
document generation should therefore be treated as validation of architectural
correctness.

The problem motivating ADR4AI is how to support the transition from industrial
requirements to candidate AI-system patterns and inspectable ADR drafts while
preserving the evidence behind each recommendation. Such support should expose
the requirement passages used to justify applicability, distinguish insufficient
information from a supported decision, and retain links between requirements,
selected patterns, and generated records. Expert review is needed to assess
whether the recommendations and their documented rationale are adequate for the
stated requirements. The intended role is traceable decision support for
architects, not autonomous acceptance of architectural decisions.

\noindent\textbf{TBD---Expert review of pattern recommendations and ADR drafts
pending Phase~3 execution and consolidation.}
The assessment of recommendation adequacy and documented rationale remains
pending; no expert-validation outcome is claimed in this Introduction.

The literature reviewed in Section~\ref{sec:background-related-work} examines
LLM-assisted architectural design, ADR decision generation, and AI-pattern
solution synthesis~\cite{cervantes2026assistants,dhar2024generating,oliveira2026designpartners}.
These lines of work motivate investigating their intersection: requirement-grounded
selection from an AI-pattern catalog linked to decision documentation and
comparison with retrieval-based candidates. ADR4AI operationalizes this
investigation through the pipeline in Figure~\ref{fig:approach-overview}.
It assesses each of \StudyPatternCount{} catalog patterns against each of
\StudyStoryCount{} industrial user stories and their acceptance criteria from
\StudySystemCount{} systems, yielding an exhaustive census of
\StudyAssessmentCount{} story--pattern pairs within this bounded sample.
The LLM classifies each pair as applicable, not applicable, or insufficient
information, requiring requirement references for positive judgments.
Each pair classified as applicable supplies the requirement context, pattern,
and applicability evidence for a traceable ADR draft with Proposed status.
A separate dense-retrieval and reranking baseline supplies a final
Top-\RagFinalCount{} candidate set for same-story overlap comparison; this is
an output-set comparison, not a test of architectural correctness or superiority.
The prepared expert-assessment protocol complements these computational stages
by asking practitioners familiar with the respective systems to assess selected
patterns and their ADRs, including decision adequacy, grounding, consequences,
perceived usefulness, and pattern applicability.
\textbf{TBD---Expert assessment execution, consolidated responses, and resulting
validation evidence pending Phase~3; this stage is planned, not reported as
completed.}
